\chapter{Quand ça rate : lire ce que Jenga dit}

Un échec de construction se lit. Ce chapitre donne l'ordre dans lequel le faire,
et les messages qu'on prend pour ce qu'ils ne sont pas.

\section{Trois étages, trois familles d'échec}

\begin{center}
\begin{tabular}{@{}lll@{}}
\toprule
\textbf{Étage} & \textbf{Le message ressemble à} & \textbf{Il faut regarder} \\
\midrule
chargement & une erreur \emph{Python} & votre \texttt{.jenga} \\
compilation & une erreur du \emph{compilateur} & votre code C++ \\
édition de liens & un \emph{symbole} manquant & vos \texttt{links} et \texttt{dependson} \\
\bottomrule
\end{tabular}
\end{center}

\begin{nkretenir}
\textbf{Identifiez l'étage AVANT de chercher.} C'est trente secondes, et cela
détermine dans quel fichier vous allez passer l'heure suivante.

\textbf{Le troisième est le plus mal compris.} Une erreur d'édition de liens
nomme un \emph{symbole} --- une fonction, une variable --- et ne nomme jamais le
module qui la contient. Elle envoie donc chercher dans le code, alors que la
cause est dans la description.
\end{nkretenir}

\section{L'échec de chargement}

\begin{center}
\begin{tabular}{@{}ll@{}}
\toprule
\textbf{Message} & \textbf{Cause la plus fréquente} \\
\midrule
\texttt{SyntaxError} & virgule, parenthèse, indentation \\
\texttt{NameError} & une fonction mal orthographiée \\
\texttt{TypeError} & une chaîne au lieu d'une liste \\
\texttt{FileNotFoundError} & un \texttt{include()} vers un fichier absent \\
\texttt{No workspace defined} & aucun \texttt{with workspace} exécuté \\
\texttt{External file not found} & un sous-module non peuplé \\
\bottomrule
\end{tabular}
\end{center}

\begin{nkretenir}
\textbf{\texttt{-{}-verbose} donne la pile complète}, et c'est le premier réflexe.
Le message dit \emph{quoi} ; la pile dit \emph{où}, y compris dans un fichier
inclus.
\end{nkretenir}

\begin{nkpiege}[La dernière ligne du tableau accuse le mauvais coupable]
\texttt{External file not found: Externals/Libs/NKGlad/NKGlad.jenga} ressemble à
une erreur de description. C'en est rarement une : \textbf{le dossier du
sous-module est vide.}

Et l'origine est plus retorse encore : dans un \emph{worktree} git,
\texttt{git submodule update -{}-init} \textbf{ne fait rien et rend 0}. Aucune
sortie, aucun message. Il faut \texttt{-{}-force}.

Tout ce qu'on lance ensuite échoue « à la construction ». \textbf{Vérifiez vos
sous-modules avant de conclure quoi que ce soit sur le code} --- sinon
l'instrument accuse le sujet à la place de son propre montage.
\end{nkpiege}

\section{L'échec de compilation}

\begin{nkretenir}
Jenga relaie le message du compilateur, en nommant le fichier :

\begin{nkcode}{Forme reelle}
Compilation Error: NkLudoJeu.cpp
  .../NkLudoJeu.h:49:42: error: non-virtual member function marked
  'override' hides virtual member function
\end{nkcode}

\textbf{Le fichier nommé en tête est celui qu'on compilait ; l'erreur peut être
dans un en-tête.} Ici, la compilation de \texttt{NkLudoJeu.cpp} échoue sur une
ligne de \texttt{NkLudoJeu.h} --- et les deux se lisent.
\end{nkretenir}

\begin{nkpiege}
\textbf{Le même en-tête peut compiler dans un fichier et échouer dans un autre.}
Cela paraît absurde et c'est courant : l'en-tête est inclus dans deux contextes
différents, avec des macros ou des \texttt{using} différents.

Mesuré sur ce dépôt avec deux types nommés \texttt{NkShaderStage} dans deux
espaces de noms : un \texttt{.cpp} compilait pendant qu'un autre échouait sur le
même en-tête.

\textbf{Dans un en-tête, qualifiez complètement.} C'est la seule parade.
\end{nkpiege}

\begin{nkretenir}
\textbf{Par défaut, la construction s'arrête au premier échec.} Il y a donc
exactement une erreur par relevé, et \texttt{Failed: 1} ne veut pas dire qu'un
seul fichier est cassé.

\texttt{-{}-keep-going} continue et donne le compte réel. À employer avant
d'estimer un chantier : le premier chiffre est presque toujours faux.
\end{nkretenir}

\section{L'échec d'édition de liens}

\begin{center}
\begin{tabular}{@{}ll@{}}
\toprule
\textbf{Symptôme} & \textbf{Ce qu'il faut vérifier} \\
\midrule
\texttt{undefined reference to <votre fonction>} & le \texttt{.cpp} est-il dans \texttt{files} ? \\
\texttt{undefined reference to WinMain} & le point d'entrée et le genre s'accordent-ils ? \\
\texttt{cannot find -lX11} & la bibliothèque est-elle installée ? \\
\texttt{framework not found} & existe-t-elle sur \emph{cette} plateforme ? \\
\texttt{multiple definition of main} & deux \texttt{main} --- voir \texttt{excludemainfiles} \\
\bottomrule
\end{tabular}
\end{center}

\begin{nkretenir}
\textbf{La première ligne est la plus fréquente, et la plus mal cherchée.} Un
symbole qu'on a pourtant écrit est introuvable parce que son fichier n'a jamais
été compilé --- motif qui ne le prend pas, ou fichier hors du dossier déclaré.

La sortie de construction porte la réponse : \texttt{Found N source file(s)}.
\textbf{Comparez N au nombre de vos fichiers} avant de chercher ailleurs.
\end{nkretenir}

\begin{nkpiege}
\textbf{Une erreur d'édition de liens ne nomme aucun de vos fichiers de
description.} \texttt{undefined reference to WinMain} ne mentionne ni votre
\texttt{.jenga}, ni le mot \texttt{windowedapp} --- et l'on cherche partout sauf
là.

\emph{Quand un message ne nomme aucun de vos fichiers, la cause est presque
toujours dans la description.}
\end{nkpiege}

\section{L'ordre de diagnostic}

\begin{center}
\begin{tabular}{@{}rll@{}}
\toprule
& \textbf{Question} & \textbf{Commande} \\
\midrule
1 & le bon workspace est-il chargé ? & \texttt{jenga info} (ligne \emph{Entry file}) \\
2 & mon projet est-il dedans ? & \texttt{jenga info} (tableau) \\
3 & mes fichiers sont-ils trouvés ? & \texttt{Found N source file(s)} \\
4 & mes drapeaux sont-ils appliqués ? & \texttt{jenga compile-flags} \\
5 & mon fichier a-t-il été compilé ? & ligne \texttt{Compiled:} \\
6 & le binaire est-il frais ? & \texttt{jenga rebuild} \\
\bottomrule
\end{tabular}
\end{center}

\begin{nkretenir}
\textbf{Cet ordre est celui du COÛT, pas celui de la probabilité.} Les quatre
premières questions se répondent sans rien construire ; la sixième coûte une
construction complète.

Quand on ne sait rien, ordonner par coût est le bon critère --- ordonner par
probabilité suppose qu'on sait déjà quelque chose.
\end{nkretenir}

\section{Les messages qui ne sont pas des erreurs}

\begin{center}
\begin{tabular}{@{}ll@{}}
\toprule
\textbf{Message} & \textbf{Ce que c'est} \\
\midrule
\texttt{No test projects found.} & normal si vous n'avez pas de tests --- mais code 1 \\
\texttt{Application console sans terminal} & une commodité de \texttt{run} \\
\texttt{APK signe avec debug.keystore} & un avertissement de \emph{destination} \\
avertissement de version embarquée & NKCode vous prévient d'une divergence \\
\bottomrule
\end{tabular}
\end{center}

\begin{nkpiege}
\textbf{Un avertissement noyé au milieu de mille lignes n'a pas été émis.}
Mesuré : l'avertissement « keystore not found » de Jenga apparaît au milieu du
journal de construction et n'est \emph{pas répété} après le \texttt{SUCCESS}
final.

Résultat : un APK sans signature, un \texttt{adb install} qui refuse, et une
enquête qui commence par le téléphone.

\emph{Quand vous écrivez un outil, répétez à la fin ce qui compte.} Un message
qui n'est dit qu'une fois, au milieu, est un message qu'on n'a pas dit.
\end{nkpiege}

\begin{nkbilan}
Trois étages d'échec --- chargement, compilation, édition de liens --- et
identifier lequel avant de chercher détermine où l'on passe l'heure suivante. Le
troisième est le plus mal compris : il nomme un symbole et jamais la description,
qui est pourtant sa cause habituelle. Un sous-module vide se présente comme une
erreur de description ; la construction s'arrête au premier échec, donc le
premier compte est faux. Diagnostiquez dans l'ordre du \emph{coût}, pas de la
probabilité --- quatre des six questions se répondent sans rien construire. Et un
avertissement dit une seule fois, au milieu d'un journal, n'a pas été dit.
\end{nkbilan}

\begin{nkquestions}
\begin{enumerate}
\item Comment reconnaît-on l'étage d'un échec ?
\item Que faire devant \texttt{External file not found} sur un \texttt{.jenga} de
      dépendance ?
\item Pourquoi \texttt{Failed: 1} est-il presque toujours faux ?
\item Quelle ligne consulter devant un \texttt{undefined reference} sur une
      fonction que vous avez écrite ?
\item Pourquoi ordonner un diagnostic par coût plutôt que par probabilité ?
\end{enumerate}
\end{nkquestions}

\begin{nkpratique}
Provoquez délibérément les cinq échecs, et \textbf{gardez le message de chacun} :

\begin{enumerate}
\item une virgule manquante dans le \texttt{.jenga} ;
\item une fonction Jenga mal orthographiée ;
\item un \texttt{files()} qui ne trouve rien ;
\item un \texttt{links} vers une bibliothèque inexistante ;
\item un \texttt{windowedapp()} sur un programme qui n'a qu'un \texttt{main}
      classique.
\end{enumerate}

Constituez-vous une page de correspondance message $\rightarrow$ cause.
\textbf{C'est le document le plus utile que vous écrirez sur Jenga}, et il ne
vaut que s'il est fait de vos propres captures.
\end{nkpratique}

\begin{nkdemo}
Prouvez qu'une erreur d'édition de liens ne désigne pas sa cause.

Retirez un fichier de votre \texttt{files()} --- sans le supprimer du disque ---
et construisez.

\textbf{Ce que la démonstration doit établir}, en trois relevés :
\begin{enumerate}
\item l'erreur nomme un \emph{symbole}, pas un fichier de description ;
\item elle ne contient aucune occurrence du mot \texttt{files} ni du nom de votre
      \texttt{.jenga} ;
\item la ligne \texttt{Found N source file(s)}, elle, a changé --- et c'était le
      seul indice, dix lignes plus haut.
\end{enumerate}

Le troisième point est la leçon : l'information était là, avant l'erreur.
\end{nkdemo}

\begin{nklogique}
Une construction qui passait hier échoue aujourd'hui, sans qu'aucun fichier
source n'ait changé.

\begin{enumerate}
\item Énumérez cinq causes possibles, dont aucune n'est dans le code C++.
\item Rangez-les par coût de vérification.
\item Deux d'entre elles ne laissent \emph{aucune trace} dans la sortie de
      construction. Lesquelles, et par quoi les prendriez-vous ?
\end{enumerate}
\end{nklogique}
